\section{Results}
\label{sec:results}

\subsection{What the matching task allows}
\label{sec:res-ladder}

We change one thing at a time, starting from the benchmark's model run in our
training loop. The benchmark itself uses matched candidates and reports positive
contributions. Here we first pair its model with \cUnm{} candidates, the usual
choice in the wider literature. This shows what the task allows; it is not a
reproduction of the benchmark's result.

With \cUnm{} candidates the benchmark's model reaches
\R{ladder_raw/L0_published/within/acc} four-way accuracy, but its contribution is
\R{ladder_raw/L0_published/within/contrib}, its decisions never change when the EEG
is shuffled (flip rate \R{ladder_raw/L0_published/within/flip}), and its
embeddings are nearly identical across windows (collapse
\R{ladder_raw/L0_published/within/collapse}). The accuracy is real, but none of it
comes from the EEG: this is the shortcut of Section~\ref{sec:failures}(ii). A
better encoder raises accuracy to \R{ladder_raw/L1_encoder/within/acc} and leaves
the contribution at \R{ladder_raw/L1_encoder/within/contrib} --- it simply reads
the audio better. A change that raises accuracy is not, by itself, evidence; only
the shuffle test shows this. The time-centred score gives the first non-zero
contribution. With matched candidates the same benchmark model scores
\R{ladder_qm/L0_published/within/acc}, which is chance.

\subsection{More streams, less credit for the brain}
\label{sec:res-hinge}

\begin{table}[t]
\caption{\textbf{More streams give more accuracy, and less credit to the brain.}
Within-listener, \cMat{}, chance $0.25$. The four fused models share
streams, data, folds, architecture and listener adaptation; they differ only in
the contribution hinge and in modality dropout. $\phi_m$ is each stream's Shapley
credit and the five add up to the contribution. The EEG-only model shows how much
credit the brain earns on its own. Appendix~\ref{app:hinge2loso} gives the same
four models on unseen listeners.}
\label{tab:hinge}
\centering\small
\setlength{\tabcolsep}{4pt}
\begin{tabular}{@{}lrrrrrrrr@{}}
\toprule
& & & & \multicolumn{5}{c}{Shapley credit $\phi_m$}\\
\cmidrule(l){5-9}
Model & Accuracy & Null & Contrib. & EEG & gaze & head & flow & fovea \\
\midrule
EEG alone
 & \R{v2onset/V4_subjadapt/within/acc} & \R{v2onset/V4_subjadapt/within/null} & \R{v2onset/V4_subjadapt/within/contrib}
 & \R{v2onset/V4_subjadapt/within/contrib} & --- & --- & --- & --- \\
\midrule
\textbf{\merit{}} (no hinge, no dropout)
 & \textbf{\R{v2hinge2/H4_nohinge_nodrop/within/acc}} & \R{v2hinge2/H4_nohinge_nodrop/within/null} & \textbf{\R{v2hinge2/H4_nohinge_nodrop/within/contrib}} & \R{v2hinge2/H4_nohinge_nodrop/within/shap/eeg} & \R{v2hinge2/H4_nohinge_nodrop/within/shap/gaze} & \R{v2hinge2/H4_nohinge_nodrop/within/shap/imu} & \R{v2hinge2/H4_nohinge_nodrop/within/shap/video} & \R{v2hinge2/H4_nohinge_nodrop/within/shap/fovea} \\
No hinge, with dropout
 & \R{v2hinge2/H2_nohinge_drop/within/acc} & \R{v2hinge2/H2_nohinge_drop/within/null} & \R{v2hinge2/H2_nohinge_drop/within/contrib} & \R{v2hinge2/H2_nohinge_drop/within/shap/eeg} & \R{v2hinge2/H2_nohinge_drop/within/shap/gaze} & \R{v2hinge2/H2_nohinge_drop/within/shap/imu} & \R{v2hinge2/H2_nohinge_drop/within/shap/video} & \R{v2hinge2/H2_nohinge_drop/within/shap/fovea} \\
With hinge, no dropout
 & \R{v2hinge2/H3_hinge_nodrop/within/acc} & \R{v2hinge2/H3_hinge_nodrop/within/null} & \R{v2hinge2/H3_hinge_nodrop/within/contrib} & \R{v2hinge2/H3_hinge_nodrop/within/shap/eeg} & \R{v2hinge2/H3_hinge_nodrop/within/shap/gaze} & \R{v2hinge2/H3_hinge_nodrop/within/shap/imu} & \R{v2hinge2/H3_hinge_nodrop/within/shap/video} & \R{v2hinge2/H3_hinge_nodrop/within/shap/fovea} \\
With hinge, with dropout
 & \R{v2hinge2/H1_hinge_drop/within/acc} & \R{v2hinge2/H1_hinge_drop/within/null} & \R{v2hinge2/H1_hinge_drop/within/contrib} & \R{v2hinge2/H1_hinge_drop/within/shap/eeg} & \R{v2hinge2/H1_hinge_drop/within/shap/gaze} & \R{v2hinge2/H1_hinge_drop/within/shap/imu} & \R{v2hinge2/H1_hinge_drop/within/shap/video} & \R{v2hinge2/H1_hinge_drop/within/shap/fovea} \\
\bottomrule
\end{tabular}
\end{table}

\paragraph{The brain's share falls by almost half.}
Adding the behavioural streams raises accuracy from \R{v2onset/V4_subjadapt/within/acc} to
\R{v2hinge2/H4_nohinge_nodrop/within/acc}. But the EEG's credit falls from \R{v2onset/V4_subjadapt/within/contrib}, when it
is the only stream, to \R{v2hinge2/H4_nohinge_nodrop/within/shap/eeg} in the fused model. More streams have
not taught the model more about the brain; they have let it rely on the brain
less. A joint shuffle cannot see this, because the joint contribution
\emph{rises}, from \R{v2onset/V4_subjadapt/within/contrib} to \R{v2hinge2/H4_nohinge_nodrop/within/contrib}. This is why the
per-stream split is the measurement, and the total is not. In the fused model the
foveal embedding earns \R{v2hinge2/H4_nohinge_nodrop/within/shap/fovea}, more than the EEG: much of the
accuracy is a statement about where the listener was looking.

\paragraph{Putting the test in the loss does not fix it.}
The contribution hinge never raised the brain's credit above the matching model
without it. With modality dropout it lowered the EEG's credit from
\R{v2hinge2/H2_nohinge_drop/within/shap/eeg} to \R{v2hinge2/H1_hinge_drop/within/shap/eeg}; without dropout it made no real
difference (\R{v2hinge2/H3_hinge_nodrop/within/shap/eeg} against \R{v2hinge2/H4_nohinge_nodrop/within/shap/eeg}). The reason is the
interaction. Dropout removes streams at random during training, so the
validation contribution that sets the hinge's weights is measured on a model that
keeps being asked to work without them, and the hinge then pushes whichever
stream looks unused --- rarely the EEG. On unseen listeners the models without the
hinge keep more credit on the brain in both pairs (Appendix~\ref{app:hinge2loso}).
A measurement that cannot be pushed up by training against it is one worth
trusting. The one training choice that did raise the brain's credit was choosing
checkpoints by contribution: in an earlier version of this experiment, at
$W=10$\,s, it moved the
EEG's credit from \R{hinge/all_nohinge_accsel/within/shap/eeg} to
\R{hinge/all_nohinge_credsel/within/shap/eeg} (Table~\ref{tab:hingew10}).

\paragraph{Which model we call \merit{}.}
\merit{} is the model without the hinge and without dropout. It is not the most
accurate of the four in either protocol. It keeps the most credit on the brain on
unseen listeners (\R{v2hinge2loso/H4_nohinge_nodrop/loso/shap/eeg}) and ties for the most within listeners
--- and credit is what we select for. Choosing it costs less than $0.02$ in
accuracy in each protocol.

\subsection{The full system, and the benchmark}
\label{sec:res-full}

The results above make one decision per window. But the listener attends to one
talker for the whole trial, so Table~\ref{tab:full} adds, one at a time: one
decision per trial, made by summing the window scores; training on random crops
of the whole trial; limiting the coupling to delays after the sound; a small
per-listener correction of the EEG channels; and a second audio channel that marks
speech onsets. The candidates stay matched and the splits stay the same.

\begin{table}[t]
\caption{\textbf{Improvement ladder}, EEG only, within-listener, \cMat{},
chance $0.25$. Each row adds one change to the row above. The last row is a
control, not a step: it limits the coupling to delays \emph{before} the sound,
where no brain response to that sound can be.}
\label{tab:full}
\centering\small
\begin{tabular}{@{}lrrr@{}}
\toprule
& Accuracy & Null & Contribution \\
\midrule
Random crops, one decision per window & \R{v2onset/V1_crops/within/acc} & \R{v2onset/V1_crops/within/null} & \R{v2onset/V1_crops/within/contrib} \\
\;+ one decision per trial         & \R{v2onset/V2_trial/within/acc} & \R{v2onset/V2_trial/within/null} & \R{v2onset/V2_trial/within/contrib} \\
\;+ only delays after the sound    & \R{v2onset/V3_causal/within/acc} & \R{v2onset/V3_causal/within/null} & \R{v2onset/V3_causal/within/contrib} \\
\textbf{\merit{}, EEG only} (+ per-listener correction) & \textbf{\R{v2onset/V4_subjadapt/within/acc}} & \R{v2onset/V4_subjadapt/within/null} & \textbf{\R{v2onset/V4_subjadapt/within/contrib}} \\
\midrule
\;\emph{control:} only delays before the sound & \R{v2onset/V3b_acausal/within/acc} & \R{v2onset/V3b_acausal/within/null} & \R{v2onset/V3b_acausal/within/contrib} \\
\bottomrule
\end{tabular}
\end{table}

\paragraph{What each change is worth.}
One decision per trial is by far the largest gain: the contribution rises from
\R{v2onset/V1_crops/within/contrib} to \R{v2onset/V2_trial/within/contrib} on the
same windows. Limiting the coupling to delays after the sound adds a little, and
the control confirms the direction: delays before the sound reach only
\R{v2onset/V3b_acausal/within/contrib}. The per-listener correction raises the
contribution from \R{v2onset/V3_causal/within/contrib} to
\R{v2onset/V4_subjadapt/within/contrib}, and the onset channel raises accuracy
from \R{v2ladder/V4_subjadapt/within/acc} to \R{v2onset/V4_subjadapt/within/acc}.

\paragraph{Against the benchmark.}
We compare with the benchmark at its longest window, $30$\,s, which --- like our
per-trial decision --- gives one decision per trial \citep{maestro2026}. There it
reports $0.591$ for EEG alone within listeners and $0.619$ on unseen listeners,
and $0.698$ and $0.674$ for its best multimodal models; its contributions are
$+0.320$ and $+0.363$ for EEG alone, and $+0.413$ for all streams on unseen
listeners. \merit{} reaches \R{v2onset/V4_subjadapt/within/acc} for EEG alone within listeners
(contribution \R{v2onset/V4_subjadapt/within/contrib}), \R{v2hinge2/H4_nohinge_nodrop/within/acc} with all streams
(\R{v2hinge2/H4_nohinge_nodrop/within/contrib}), and \R{v2hinge2loso/H4_nohinge_nodrop/loso/acc} with all streams on unseen
listeners (\R{v2hinge2loso/H4_nohinge_nodrop/loso/contrib}, positive in \R{v2hinge2loso/H4_nohinge_nodrop/loso/pos} folds). Three of
the four comparisons improve on both accuracy and contribution.

The fourth, EEG alone on unseen listeners, does not: \R{v2bestloso/eeg/loso/acc}
against $0.619$, with contribution \R{v2bestloso/eeg/loso/contrib}. The per-listener
correction behind our largest within-listener gain cannot, by design, help a
listener it has never seen. Our gains on the \emph{brain} signal are therefore
shown within listeners.

\subsection{Is the EEG signal real?}
\label{sec:res-mm}

Four checks, with details in Appendix~\ref{app:more}.
\textbf{Same-talker candidates.} When all candidates come from the attended
talker's own speech, the audio alone is at chance and no orienting cue can help;
EEG alone still contributes \R{mm3/eeg/within/contrib} at $K\!=\!3$ and
\R{mm2/eeg/within/contrib} at $K\!=\!2$.
\textbf{Direction in time.} Limited to delays of $+125$ to $+1109$\,ms after the
sound, where a brain response must be, the contribution is
\R{lag_qm/causal/within/contrib}; the mirror image before the sound gives
\R{lag_qm/acausal/within/contrib}. This rules out an electrical echo of the sound
at zero delay, and slow drift that is the same in both directions.
\textbf{Window length.} The contribution grows with the window, as it should for
this task \citep{geirnaert2021} (Figure~\ref{fig:window}). At short windows the
behavioural streams take over: at $2$\,s the EEG alone contributes
\R{w2/eeg/within/contrib} while the foveal embedding alone contributes
\R{w2/fovea/within/contrib} (Appendix~\ref{app:shortwin}).
\textbf{Other datasets.} On two public two-talker datasets the audio probe is at
chance on DTU but \R{extprobe/kul/w10/excess} above chance on KU~Leuven,
because one of its two talkers is the attended one in \R{extprior/kul/nrole} of \R{extprior/kul/ntrials} trials
(Appendix~\ref{app:external}). On KU~Leuven and DTU the coupling branch reaches
\R{extkulwithinst/eeg_centred/within/acc} and \R{extdtuwithinst/eeg_centred/within/acc}
even with same-talker candidates, but the direction-in-time check does not pass:
delays before and after the sound score the same. We therefore do not count these
as brain decoding (Appendix~\ref{app:externalcoupling}). At the $5$\,s window used by recent work,
the coupling branch is more accurate on unseen listeners of both datasets than
NeuroToken \citep{alavi2026neurotoken}, the strongest prior result without data
leakage (Table~\ref{tab:published}). A plain accuracy report
would have done so.
